% Bilan énergétique de la réaction 2 H2 + O2 -> 2 H2O (gaz) à partir des énergies de liaison (corrigé de l'exercice 5, chapitre 10)
\begin{tikzpicture}[x=1cm,y=0.0021cm]
\draw[otMuted,line width=0.5pt,-{Stealth[length=2mm]}] (-2.0,-650) -- (-2.0,1650) node[above,pcLegende] {énergie};
\draw[otNavy,line width=1.2pt] (0,0) -- (2.4,0) node[right,pcLegende,text=otNavy] {$\ce{2H2 + O2}$ (gaz)};
\draw[otMuted,line width=1.2pt] (2.6,1370) -- (6.2,1370) node[right,pcLegende,align=left] {atomes séparés :\\4 H $+$ 2 O};
\draw[otGreenDark,line width=1.2pt] (6.4,-482) -- (8.8,-482) node[right,pcLegende,text=otGreenDark] {$\ce{2H2O}$ (gaz)};
\draw[otRed!80!black,line width=1pt,-{Stealth[length=2.5mm]}] (1.6,40) -- (3.4,1330)
  node[midway,left,pcLegende,text=otRed!80!black,align=right] {rupture de 2 H--H et 1 O=O\\$+\qty{1370}{kJ}$ reçus};
\draw[otGreenDark,line width=1pt,-{Stealth[length=2.5mm]}] (5.4,1330) -- (7.4,-440)
  node[pos=0.42,right,pcLegende,text=otGreenDark,align=left] {formation de 4 O--H\\$\qty{1852}{kJ}$ libérés};
\draw[otOrange!85!black,line width=0.9pt,dash pattern=on 3pt off 1.5pt,-{Stealth[length=2.3mm]}] (2.4,0) -- (6.4,-482)
  node[pos=0.45,below left,pcLegende,text=otOrange!75!black] {$E_\mathrm{r}=1370-1852=\qty{-482}{kJ.mol^{-1}}$};
\end{tikzpicture}
\par\smallskip\legende{Pour l'équation écrite, qui forme deux moles d'eau ; par mole de dihydrogène, l'énergie libérée vaut \qty{241}{kJ}.}
